\paragraph{Pauli-Y Gate}\label{subsubsection: Pauli-Y Gate}

The \definitionWithSpecificIndex{Pauli-Y Gate}{Single Qubit Gates: Pauli-Y Gate}{}, denoted $Y$, is one of the three fundamental \textbf{Pauli matrices} in quantum mechanics. It is less intuitive than Pauli-X and Z, but equally important, especially for its role in complex rotations and quantum interference.

\highspace
\begin{flushleft}
    \textcolor{Green3}{\faIcon{square-root-alt} \textbf{Matrix Representation}}
\end{flushleft}
\begin{equation*}
    Y = \sigma_{Y} = \begin{pmatrix}
        0 & -i \\ i & 0
    \end{pmatrix}
\end{equation*}

\highspace
\begin{flushleft}
    \textcolor{Green3}{\faIcon{book} \textbf{Action on a General Qubit}}
\end{flushleft}
At first glance, this looks similar to the Pauli-X gate because it exchanges $\ket{0}$ and $\ket{1}$. But there is something extra: $X$ swaps the two basis states (i.e., $\ket{0} \leftrightarrow \ket{1}$), while $Y$ swaps them \textbf{and} multiplies them by a phases of $\pm i$. So Pauli-Y is not just a bit flip (like Pauli-X), it is a bit flip together with a phase modification.

\highspace
Given a qubit in state:
\begin{equation*}
	\ket{\psi} = a\ket{0} + b\ket{1}
\end{equation*}
Apply the Y gate:
\begin{equation*}
	Y\ket{\psi} = aY\ket{0} + bY\ket{1} = a\left(-i \ket{1}\right) + b\left(i \ket{0}\right) = ib\ket{0} - ia\ket{1}
\end{equation*}
So the \textbf{amplitudes are exchanged} (as with Pauli-X), but each one also acquires a \textbf{phase factor of $i$ or $-i$}.

\highspace
\begin{flushleft}
    \textcolor{Green3}{\faIcon{question-circle} \textbf{Geometric Interpretation: Rotation on the Bloch Sphere}}
\end{flushleft}
The Pauli-Y gate corresponds to a \textbf{rotation around the $y$-axis by $\pi$ radians ($180 \degree$)}. The effect on Bloch sphere coordinates:
\begin{itemize}
    \item $\theta \rightarrow \pi - \theta$
    \item $\psi \rightarrow \pi - \psi$
\end{itemize}
This rotates the qubit across the $y$-axis, \textbf{altering both amplitudes and their phases}.

\highspace
\begin{flushleft}
    \textcolor{Green3}{\faIcon{book} \textbf{Pauli-Y as a Composite Gate}}
\end{flushleft}
A very useful identity is:
\begin{equation}
    Y = i X Z
\end{equation}
This means that Pauli-Y can be understood as the combination of:
\begin{enumerate}
    \item A Pauli-Z gate (phase flip);
    \item Followed by a Pauli-X gate (bit flip);
    \item Up to a global phase factor of $i$.
\end{enumerate}
\begin{equation*}
    XZ = \begin{pmatrix}
        0 & 1 \\ 1 & 0
    \end{pmatrix}
    \begin{pmatrix}
        1 & 0 \\ 0 & -1
    \end{pmatrix}
    =
    \begin{pmatrix}
        0 & -1 \\ 1 & 0
    \end{pmatrix}
\end{equation*}
Now, multiply by $i$:
\begin{equation*}
    iXZ = i \begin{pmatrix}
        0 & -1 \\ 1 & 0
    \end{pmatrix}
    =
    \begin{pmatrix}
        0 & -i \\ i & 0
    \end{pmatrix}
    = Y
\end{equation*}
This identity confirms that \textbf{Pauli-Y can be viewed as a combination of Pauli-X and Pauli-Z with a global phase factor $i$}, which does not affect measurement outcomes.

\begin{figure}[!htp]
    \centering
    \tdplotsetmaincoords{70}{120}
    \begin{tikzpicture}[tdplot_main_coords, scale=3.4, >=Latex]

        \def\r{1}

        % Axes
        \draw[->, thick] (0,0,0) -- (1.25,0,0) node[anchor=north east] {$x$};
        \draw[->, very thick, Green3] (0,-1.25,0) -- (0,1.35,0) node[anchor=west] {$y$};
        \draw[->, thick] (0,0,0) -- (0,0,1.25) node[anchor=south] {$z$};

        % Negative axis hints
        \draw[dashed, thin, black!45] (-1.10,0,0) -- (0,0,0);
        \draw[dashed, thin, black!45] (0,0,-1.10) -- (0,0,0);

        % Bloch sphere
        \tdplotdrawarc[dashed]{(0,0,0)}{\r}{0}{180}{}{}
        \tdplotdrawarc[thick]{(0,0,0)}{\r}{180}{360}{}{}

        \tdplotsetrotatedcoords{90}{0}{90}
        \tdplotdrawarc[tdplot_rotated_coords,dashed]{(0,0,0)}{\r}{0}{180}{}{}
        \tdplotdrawarc[tdplot_rotated_coords,thick]{(0,0,0)}{\r}{180}{360}{}{}
        \tdplotsetrotatedcoords{0}{0}{0}

        % Eigenstates on y-axis
        \fill[Green3] (0,1,0) circle (0.018);
        \node[Green3, anchor=west] at (0,1.08,-0.1) {$\ket{i}$};

        \fill[Green3] (0,-1,0) circle (0.018);
        \node[Green3, anchor=east] at (0,-1.10,0.1) {$\ket{-i}$};

        % Computational basis states for reference
        \node[anchor=south east] at (0,0,1.05) {$\ket{0}$};
        \node[anchor=north east] at (0,0,-1.05) {$\ket{1}$};

        % Manual magenta rotation arrow around the y-axis
        % The curve is drawn in 3D coordinates to make the rotation visually clear.
        \draw[->, very thick, magenta]
            (0.45,0.75,0.35)
            .. controls (-0.45,0.55,0.65) and (-0.55,0.45,-0.15) ..
            (0.00,0.55,-0.70);

        \node[magenta, anchor=west] at (0.48,0.85,0.42) {$\pi$ rotation around $y$};
    \end{tikzpicture}
    \caption{Eigenstates of the Pauli-$Y$ gate. Since $\ket{i}$ and $\ket{-i}$ lie on the $y$-axis, they are not moved by a rotation around that axis. The state $\ket{-i}$ only acquires a global phase $-1$.}
    \label{fig:pauli-y-eigenstates}
\end{figure}

\highspace
\begin{flushleft}
    \textcolor{Green3}{\faIcon{book} \textbf{Eigenstates of Y Gate}}
\end{flushleft}
To find the eigenstates, we solve:
\begin{equation*}
    Y\ket{v} = \lambda \ket{v}
\end{equation*}
Where $\lambda$ is the eigenvalue. The result is:
\begin{equation}
    \ket{i} = \dfrac{1}{\sqrt{2}} \left(\ket{0} + i\ket{1}\right)
\end{equation}
\begin{equation}
    \ket{-i} = \dfrac{1}{\sqrt{2}} \left(\ket{0} - i\ket{1}\right)
\end{equation}
With eigenvalues:
\begin{equation}
    Y\ket{i} = +1 \ket{i} \qquad Y\ket{-i} = -1 \ket{-i}
\end{equation}

\highspace
\begin{flushleft}
    \textcolor{Green3}{\faIcon{list-ul} \textbf{Key Properties of Y Gate}}
\end{flushleft}
\begin{enumerate}
    \item \textbf{Unitary}: $Y^\dagger Y = I$. So it is a valid quantum gate that preserves normalization.
    \item \textbf{Hermitian}: $Y^\dagger = Y$. So it is also an observable, meaning it corresponds to a measurable quantity in quantum mechanics.
    \item \textbf{Self-inverse}: $Y^2 = I$. Applying the Y gate twice returns the qubit to its original state.
\end{enumerate}
